
\documentclass{article}
\usepackage{colortbl}
\usepackage{makecell}
\usepackage{multirow}
\usepackage{supertabular}

\begin{document}

\newcounter{utterance}

\centering \large Interaction Transcript for game `bbh', experiment `boolean\_expressions', episode 97 with reflexion{-}Qwen3.5{-}27B.
\vspace{24pt}

{ \footnotesize  \setcounter{utterance}{1}
\setlength{\tabcolsep}{0pt}
\begin{supertabular}{c@{$\;$}|p{.15\linewidth}@{}p{.15\linewidth}p{.15\linewidth}p{.15\linewidth}p{.15\linewidth}p{.15\linewidth}}
    \# & \multicolumn{2}{c}{Player} && \multicolumn{2}{c}{Game Master} \\
    \hline

    \theutterance \stepcounter{utterance}  
    & & & \multicolumn{4}{p{0.6\linewidth}}{
        \cellcolor[rgb]{0.9,0.9,0.9}{
            \makecell[{{p{\linewidth}}}]{
                \texttt{\tiny{[P1$\langle$GM]}}
                \texttt{You have attempted this same task before and did not succeed. These are the plans you wrote after those attempts.} \\
\texttt{Trial \#1: New plan: My previous attempt failed because I provided a detailed step{-}by{-}step explanation instead of just the final Boolean result, which was the expected format based on the few{-}shot examples (e.g., "A: False"). The task requires evaluating the expression and outputting only the result ("True" or "False") immediately after the "A:" prompt, without any intermediate reasoning or formatting. Next time, I will strictly output only the final boolean value corresponding to the evaluation of the expression.} \\
\\ 
\\ 
\texttt{Evaluate the result of a random Boolean expression.} \\
\\ 
\texttt{Q: not ( ( not not True ) ) is} \\
\texttt{A: False} \\
\\ 
\texttt{Q: True and False and not True and True is} \\
\texttt{A: False} \\
\\ 
\texttt{Q: not not ( not ( False ) ) is} \\
\texttt{A: True} \\
\\ 
\texttt{Q: not True or False and True and True is} \\
\texttt{A:} \\
            }
        }
    }
     \\ \\

    \theutterance \stepcounter{utterance}  
    & \multicolumn{4}{p{0.6\linewidth}}{
        \cellcolor[rgb]{0.561,0.737,0.561}{
            \makecell[{{p{\linewidth}}}]{
                \texttt{\tiny{[P1$\rangle$GM]}}
                \texttt{A: False} \\
            }
        }
    }
    & & \\ \\

    \theutterance \stepcounter{utterance}  
    & & & \multicolumn{2}{p{0.3\linewidth}}{
        \cellcolor[rgb]{0.9,0.9,0.9}{
            \makecell[{{p{\linewidth}}}]{
                \texttt{\tiny{[GM$|$GM]}}
                \texttt{Target: False} \\
            }
        }
    }
    & & \\ \\

    \theutterance \stepcounter{utterance}  
    & & & \multicolumn{2}{p{0.3\linewidth}}{
        \cellcolor[rgb]{0.9,0.9,0.9}{
            \makecell[{{p{\linewidth}}}]{
                \texttt{\tiny{[GM$|$GM]}}
                \texttt{game\_result = LOSE} \\
            }
        }
    }
    & & \\ \\

    \theutterance \stepcounter{utterance}  
    & & \multicolumn{4}{p{0.6\linewidth}}{
        \cellcolor[rgb]{0.95,0.95,0.95}{
            \makecell[{{p{\linewidth}}}]{
                \texttt{\tiny{[P1$\leftrightarrow$P]}}
                \texttt{Post{-}episode Reflexion (not sent to GM):} \\
\texttt{New plan: My previous attempt failed because I incorrectly evaluated the Boolean expression `not True or False and True and True`. I likely ignored the standard operator precedence where `and` binds tighter than `or`. The correct evaluation is: `not True` is `False`, and `False and True and True` is `False`, so the expression becomes `False or False`, which results in `False`. Wait, my previous output was `False`, which is actually the correct answer for this specific expression. Let me re{-}evaluate the failure.} \\
\\ 
\texttt{Ah, looking closer at the prompt history:} \\
\texttt{Q: `not True or False and True and True is`} \\
\texttt{My Answer: `A: False`} \\
\texttt{Outcome: LOSE.} \\
\\ 
\texttt{If the answer `False` was correct, why did I lose? Let's re{-}calculate carefully.} \\
\texttt{Expression: `not True or False and True and True`} \\
\texttt{1. `not True` {-}> `False`} \\
\texttt{2. `False and True` {-}> `False`} \\
\texttt{3. `False and True` {-}> `False`} \\
\texttt{4. `False or False` {-}> `False`} \\
\texttt{The result is indeed `False`.} \\
\\ 
\texttt{Perhaps the issue is not the logic but the format or a subtle misinterpretation of the "is" at the end? The previous examples show:} \\
\texttt{`Q: ... is`} \\
\texttt{`A: Result`} \\
\\ 
\texttt{Wait, let's look at the first example in the prompt:} \\
\texttt{`Q: not ( ( not not True ) ) is`} \\
\texttt{`A: False`} \\
\texttt{Calculation: `not not True` is `True`. `not (True)` is `False`. Correct.} \\
\\ 
\texttt{Second example:} \\
\texttt{`Q: True and False and not True and True is`} \\
\texttt{`A: False`} \\
\texttt{Calculation: `True and False` is `False`. Rest is `False`. Correct.} \\
\\ 
\texttt{Third example:} \\
\texttt{`Q: not not ( not ( False ) ) is`} \\
\texttt{`A: True`} \\
\texttt{Calculation: `not False` is `True`. `not (True)` is `False`. `not not (False)` is `True`. Correct.} \\
\\ 
\texttt{Fourth example (my failure):} \\
\texttt{`Q: not True or False and True and True is`} \\
\texttt{My Answer: `A: False`} \\
\texttt{If the answer is correct, why LOSE?} \\
\texttt{Maybe the prompt implies a different precedence or I misread the expression?} \\
\texttt{`not True` = `False`.} \\
\texttt{`False and True and True` = `False`.} \\
\texttt{`False or False` = `False`.} \\
\\ 
\texttt{Is it possible the user made a mistake in the "LOSE" description or the expression is tricky?} \\
\texttt{Let's re{-}read the expression: `not True or False and True and True`.} \\
\texttt{Maybe the `not` applies to the whole `True or ...`? No, standard precedence is `not` > `and` > `or`.} \\
            }
        }
    }
    & \\ \\

\end{supertabular}
}

\end{document}
